\usepackage{amsmath,amssymb,amsthm}
\usepackage[margin=1in]{geometry}
\usepackage{xcolor}
%\usepackage{url}
%\usepackage{float}
%\usepackage{array}
\usepackage{comment}
\usepackage{booktabs}
\usepackage{graphicx}
\usepackage{tcolorbox}
%\usepackage{tikz}
%\usetikzlibrary{decorations.pathreplacing,arrows.meta,calc,patterns}
\usepackage{enumitem}
%\usepackage{csvsimple}
\usepackage{hyperref}
\usepackage[capitalize,sort]{cleveref}
%\usepackage{algorithm,algpseudocode}

% colors
\colorlet{textColor}{black}
\colorlet{bgColor}{white}
\colorlet{textRed}{red!50!black}
\colorlet{textGreen}{green!50!black}
\colorlet{textBlue}{blue!50!black}
\colorlet{textHeavy}{black}
\colorlet{dimColor}{textColor!50!bgColor}

\hypersetup{
colorlinks,linkcolor=textRed,citecolor=textRed,urlcolor=textBlue,
hypertexnames=false,bookmarksnumbered=true,final
}

% macros
\newcommand*{\Th}{^{\textrm{th}}}
\newcommand*{\abs}[1]{\lvert #1 \rvert}
\newcommand*{\defeq}{:=}
\DeclareMathOperator*{\E}{\mathbb{E}}
\DeclareMathOperator*{\Var}{Var}
\DeclareMathOperator*{\argmin}{argmin}
\DeclareMathOperator*{\argmax}{argmax}
\renewcommand*{\Pr}{\operatorname{\mathrm{P}}}

% theorems, lemmas, definitions
\theoremstyle{plain}
\newtheorem{theorem}{Theorem}
\newtheorem{lemma}[theorem]{Lemma}
\theoremstyle{definition}
\newtheorem{definition}{Definition}
\newtheorem{example}{Example}
\newtheorem{remark}{Remark}

% wrap URLs
\Urlmuskip=0mu plus 0.1mu
\makeatletter
\g@addto@macro{\UrlBreaks}{%
\do\/%
\do\a\do\b\do\c\do\d\do\e\do\f\do\g\do\h\do\i\do\j\do\k\do\l\do\m%
\do\n\do\o\do\p\do\q\do\r\do\s\do\t\do\u\do\v\do\w\do\x\do\y\do\z%
\do\A\do\B\do\C\do\D\do\E\do\F\do\G\do\H\do\I\do\J\do\K\do\L\do\M%
\do\N\do\O\do\P\do\Q\do\R\do\S\do\T\do\U\do\V\do\W\do\X\do\Y\do\Z%
\do\0\do\1\do\2\do\3\do\4\do\5\do\6\do\7\do\8\do\9%
}
\makeatother

% tighter lists
\newenvironment*{tightemize}{\begin{itemize}[noitemsep]}{\end{itemize}}%
\newenvironment*{tightenum}{\begin{enumerate}[noitemsep]}{\end{enumerate}}%

% customizations
\allowdisplaybreaks
%\defaultaddspace=0.0em  % spacing in tables

\newcommand*{\bibPrep}{
\phantomsection
\addcontentsline{toc}{section}{References}
\bibliographystyle{plainurl}
}


\newcommand*{\nhat}{\widehat{n}}
\newcommand*{\textCP}[2]{\Pr(\text{#1} \mid #2)}

\title{IE 300 Case Study 2: Predicting Flight Delays}
\author{\empty}
\date{\empty}

\begin{document}

\maketitle
%\setlength{\parindent}{0pt}
\setlength{\parskip}{0.1em}

\section{Objectives}

The purpose of this study is to develop a very popular predictive model called Naive-Bayes.
You will use Python as a tool to analyze data of flight delays in order to determine the likelihood of delays.
Through this case study, you will learn how to:
\begin{tightenum}
\item Read and analyze data using Python.
\item Apply the concepts of conditional probability and Bayes' theorem.
\item Develop a simple Naive-Bayes classifier model with $m$-estimates.
\end{tightenum}

\section{Predicting Delays}

As any frequent flier knows, flight delays are an inevitable part of commercial air travel.
In the previous case study, you were able to determine which airline is the fastest for a particular route.
In doing so, you had to take into account departure and arrival delays of each airline.
While the cause for delays are largely unpredictable, we can get some insights on
which flights are more likely to be delayed without taking weather conditions into consideration.

An obvious predictor of delays is the airline itself.
For example, in 2014, only 74.34\% of American Airline’s flights were on-time
while Delta had a 85.21\% on-time record \cite{skift2014worlds}.
How can we use historical data to make predictions about delays?
You may have already learned about several predictive models such as ordinary least-squares.
We will study a very popular model in statistics and machine learning called Naive-Bayes.

\subsection{Naive Bayes Classifier}

Recall that Bayes' theorem tells you the probability of an event based on
conditions possibly related to that event. That is, given events $A$ and $B$, we have
\begin{equation}
\Pr(A \mid B) = \frac{\Pr(A)\Pr(B \mid A)}{\Pr(B)}.
\end{equation}

The Naive-Bayes classifier is based on Bayes' theorem with the assumption of independence between attributes.
Suppose we have a vector $X = (x_1, x_2, \ldots, x_n)$ representing $n$ \textbf{attributes}
(independent variables). The Naive-Bayes' classifier helps us compute $\Pr(C_k \mid X)$,
where $C_1$, $C_2$, $\ldots$, represent possible outcomes (\textbf{classes}).

If we are trying to decide whether a particular flight is likely to be delayed,
the outcome would be from the set $\{Y, N\}$, corresponding to the
flight being delayed ($Y$) or not ($N$).
The vector $X$ would contain information about the flight
(e.g., airline, origin airport, destination airport)
that could give us some insight into which class ($Y$ or $N$) will occur.
One approach to predicting the class is to pick the value $C_k$ that is most probable,
given that we know the values of the attributes.
In other words, we pick the $C_k$ that maximizes $\Pr(C_k \mid X)$.
Since the vector $X$ is known, and $\Pr(C_k \mid X) = \Pr(C_k \cap X)/\Pr(X)$
by the definition of conditional probability,
we get that $\Pr(X)$ will be the same for all possible classes,
and we can instead focus on selecting the class that maximizes $\Pr(C_k \cap X)$.
We call $\Pr(C_k \cap X)$ the \textbf{score} for class $C_k$,
and we call $\Pr(C_k \mid X)$ the \textbf{posterior probability} for class $C_k$.

Noting that the attribute vector $X$ typically contains values of several attributes,
we can rewrite $\Pr(C_k \cap X)$ as $\Pr(C_k \cap x_1 \cap x_2 \cap \ldots \cap x_n)$,
where the event $x_i$ tells us the value of attribute $i$.
In practice, the symbol $\cap$ is often replaced with a comma when intersecting many events.
By repeatedly applying the definition of conditional probability, we get
\begin{align}
\Pr(C_k \cap X) &= \Pr(C_k)\Pr(X \mid C_k) \nonumber
\\ &= \Pr(C_k)\Pr(x_1 \cap \ldots \cap x_n \mid C_k) \nonumber
\\ &= \Pr(C_k)\Pr(x_1 \mid C_k)\Pr(x_2 \cap \ldots \cap x_n \mid C_k) \nonumber
\\ &\quad\vdots \nonumber
\\ &= \Pr(C_k)\Pr(x_1 \mid C_k)\Pr(x_2 \mid C_k \cap x_1)
    \ldots\Pr(x_n \mid C_k \cap x_1 \cap ... \cap x_{n-1}).
    \label{eq:prob-prod}
\end{align}

This decomposes $\Pr(C_k \cap X)$ into a product of the probability of outcome $C_k$
and conditional probabilities of each attribute.
We assume conditional independence between the attributes (hence, the ``naive'' portion of the name).
That is,
\begin{equation}
\Pr\bigg(x_i \Bigm\vert C_k \cap \bigg(\bigcap_{j \neq i} x_j\bigg)\bigg) = \Pr(x_i \mid C_k)
\end{equation}
for all $i \in \{1, 2, \ldots, n\}$.
This allows us to discard all conditions except $C_k$ in equation \eqref{eq:prob-prod}.
Hence, the equation simplifies to
\begin{equation}
\label{eq:naive-prod}
\Pr(C_k \cap X) = \Pr(C_k)\prod_{i=1}^n \Pr(x_i \mid C_k).
\end{equation}

With this definition, we can define a \textbf{classifier} (a decision rule).
A common classifier is to pick the most probable class, as discussed above.
This is called the \textbf{maximum a posterior} (MAP) rule, which assigns the label
\begin{equation}
\argmax_k \;\Pr(C_k)\prod_{i=1}^n \Pr(x_i \mid C_k).
\end{equation}

\subsection{\texorpdfstring{\boldmath $m$}{m}-Estimates}

The conditional densities $\Pr(x_i \mid C_k)$ can be defined in different ways (e.g., Normal, Poisson)
depending on the problem. If our dataset is small, we often do not know the underlying distributions.
When the attribute $x_i$ is discrete, an intuitive way to estimate $\Pr(x_i \mid C_k)$ is to define
\begin{equation}
\Pr(x_i \mid C_k) = \frac{\nhat}{n}.
\end{equation}
Here $\nhat$ is the number of observations where the class is $C_k$ and the $i\Th$ attribute is $x_i$,
and $n$ is the number of observations where the class is $C_k$.
However, for small datasets, we may find that $\nhat = 0$ or $n = 0$, which is a problem.
A way to avoid this problem is to use $m$-estimates.
\begin{equation}
\Pr(x_i \mid C_k) = \frac{\nhat + mp}{n + m}.
\end{equation}
Here $m$ is called the \emph{equivalent sample size} and $p$ is the \emph{prior estimate} of $\Pr(x_i \mid C_k)$.
A typical choice for $p$ is $1/(\text{no.\ of possible values for attribute } i)$, which assumes that
all attribute values are equally likely.
The idea behind $m$-estimates is to pretend we have $m$ extra observations with
class $C_k$, with $m \cdot p$ of them having the $i\Th$ attribute $x_i$.
Essentially, $m$ says how confident we are of our prior estimate $p$,
since as $m$ increases, we have $\Pr(x_i \mid C_k) \to p$.

Thus, we can compute $\Pr(x_i \mid C_k)$ for each attribute $x_i$,
and then we can compute the score of class $C_k$ using equation \eqref{eq:naive-prod}.

\subsection{Posterior Probabilities}

To compute the posterior probability $\Pr(C_k \mid X)$, we use the fact that
posterior probabilities of all classes must sum to 1, i.e.,
\begin{equation}
\sum_k \Pr(C_k \mid X) = 1.
\end{equation}
Since $\Pr(C_k \mid X) = \Pr(C_k \cap X)/\Pr(X)$ for each class $C_k$, we get
\begin{equation}
\Pr(X) = \sum_k \Pr(C_k \cap X).
\end{equation}

Hence, one can compute posterior probabilities as follows:
\begin{tightenum}
\item For each class $C_k$, compute the score $\Pr(C_k \cap X)$.
\item Sum the scores of all classes to obtain $\Pr(X)$.
\item Divide the score of each class $C_k$ by $\Pr(X)$ to obtain the posterior probability $\Pr(C_k \mid X)$.
\end{tightenum}

\section{Example 1}

Consider the following table (also available as \texttt{flight-delay.example.csv})
containing information of delays.

\begin{center}
\begin{tabular}{ccccc}
\toprule
\# & Origin & Destination & Airline & Delay
\\ \midrule
1 & DFW & ORD & Delta & Y \\
2 & DFW & ORD & Delta & N \\
3 & DFW & ORD & Delta & Y \\
4 & LAX & ORD & Delta & N \\
5 & LAX & ORD & United & Y \\
6 & LAX & LGA & United & N \\
7 & LAX & LGA & United & Y \\
8 & LAX & LGA & Delta & N \\
9 & DFW & LGA & United & N \\
10& DFW & ORD & United & Y \\
\bottomrule
\end{tabular}
\end{center}

The class is the variable Delay and the attributes are Origin, Destination, and Airline.
Classify DFW-LGA on Delta using Naive-Bayes with MAP and $m$-estimates. Assume $m = 3$.

First, note that DFW-LGA on Delta is not contained in the table.
However, we can compute the delay scores for this flight using equation \eqref{eq:naive-prod}.
Using an equivalent sample size of $m = 3$, we compute the relevant conditional probabilities as follows:
\begin{align*}
\textCP{DFW}{Y} &= \frac{3 + 3/2}{5 + 3} = \frac{9}{16} = 0.5625
& \textCP{DFW}{N} &= \frac{2 + 3/2}{5 + 3} = \frac{7}{16} = 0.4375
\\ \textCP{LGA}{Y} &= \frac{1 + 3/2}{5 + 3} = \frac{5}{16} = 0.3125
& \textCP{LGA}{N} &= \frac{3 + 3/2}{5 + 3} = \frac{9}{16} = 0.5625
\\ \textCP{Delta}{Y} &= \frac{2 + 3/2}{5 + 3} = \frac{7}{16} = 0.4375
& \textCP{Delta}{N} &= \frac{3 + 3/2}{5 + 3} = \frac{9}{16} = 0.5625
\end{align*}

For $\textCP{DFW}{Y}$, we have 5 observations where $\text{Delay} = Y$,
of which three observations have $x_i = \text{DFW}$. Thus, $n = 5$ and $\nhat = 3$.
Origin takes two possible values (DFW and LAX), so $p = 1/2$.
Since each of $Y$ and $N$ occur in 5 of the 10 observations, we have $\Pr(Y) = \Pr(N) = 0.5$.
On applying equation \eqref{eq:naive-prod}, we get the scores of $Y$ and $N$:
\begin{align*}
\Pr(Y \cap (\text{DFW, LGA, Delta})) &= \Pr(Y)\textCP{DFW}{Y}\textCP{LGA}{Y}\textCP{Delta}{Y}
    \\ &= \frac{1}{2}\cdot\frac{9}{16}\cdot\frac{5}{16}\cdot\frac{7}{16}
    = \frac{315}{8192} \approx 0.0385.
\\ \Pr(N \cap (\text{DFW, LGA, Delta})) &= \Pr(N)\textCP{DFW}{N}\textCP{LGA}{N}\textCP{Delta}{N}
    \\ &= \frac{1}{2}\cdot\frac{7}{16}\cdot\frac{9}{16}\cdot\frac{9}{16}
    = \frac{567}{8192} \approx 0.0692.
\end{align*}
Since $N$ has a larger score, we classify DFW-LGA on Delta as $N$,
i.e., we predict that the flight from DFW to LGA on Delta is likely to not be delayed.

Moreover, the posterior probabilities are:
\begin{align*}
\Pr(Y \mid \text{DFW, LGA, Delta}) &= \frac{315}{315 + 567} = \frac{5}{14} \approx 0.357,
\\ \Pr(N \mid \text{DFW, LGA, Delta}) &= \frac{567}{315 + 567} = \frac{9}{14} \approx 0.643.
\end{align*}

\section{Exercise 1}

Using the following table (also available as \texttt{flight-delay.small.csv}),
will the flight SEA-ATL on Southwest with good weather be delayed?
Use $m$-estimates for the conditional probabilities with $m = 4$.

\begin{center}
\begin{tabular}{cccccc}
\toprule
\# & Origin & Destination & Airline & Weather & Delay
\\ \midrule
1 & SEA & ATL & WN & Poor & N \\
2 & SEA & BOS & AA & Good & Y \\
3 & SEA & ATL & UN & Poor & Y \\
4 & SFO & BOS & WN & Poor & Y \\
5 & SFO & ATL & AA & Good & N \\
6 & SFO & BOS & UN & Poor & Y \\
7 & SFO & BOS & UN & Good & N \\
8 & SEA & BOS & WN & Poor & Y \\
\bottomrule
\end{tabular}
\end{center}

\section{Real World Data}

Let us now use real historical flight data to predict flight delays.
The dataset \texttt{flight-delay.large.csv} contains 2,346 observations of flights
departing from JFK or SFO to ATL, LAS or ORD in November 2015.
There are 5 variables.

\begin{center}
\begin{tabular}{ll}
\toprule
Variable Name & Description
\\ \midrule
Carrier & IATA Carrier identification \\
Origin & Origin airport \\
Destination & Destination airport \\
Depature Delay & Difference between scheduled and actual departure time (minutes) \\
Arrival Delay & Difference between scheduled and actual arrival time (minutes) \\
\bottomrule
\end{tabular}
\end{center}

\subsection{Exercise 2}

Using \texttt{flight-delay.large.csv}, write Python code to perform the following:

\begin{enumerate}[label=(\alph*)]
\item Read the dataset.
\item Compute the total delay time for each flight.
    The total delay time is the sum of departure delay and arrival delay time.
\item Assign any flight with total delay time greater than 15 minutes with $Y$ for being delayed.
    Otherwise, the flight is not delayed and is assigned $N$.
\end{enumerate}

\subsection{Exercise 3}

Using your Python code from the previous exercise, for each of the following flights,
compute the probabilities of delay and not delay given the route and carrier
and then classify whether the flight is delayed using Naive-Bayes (assume $m = 3$).

\begin{enumerate}[noitemsep,label=(\alph*)]
\item JFK-LAS on American Airlines (AA).
\item JFK-LAS on JetBlue (B6).
\item SFO-ORD on Virgin Airlines (VX).
\item SFO-ORD on Southwest Airlines (WN).
\end{enumerate}

Specifically, your output should contain:
\begin{enumerate}[noitemsep]
\item $\Pr(\text{Delay} \mid \text{Origin, Destination, Carrier})$.
\item $\Pr(\text{Not delay} \mid \text{Origin, Destination, Carrier})$.
\item Classification of whether the flight is delayed or not.
\end{enumerate}

\bibPrep{}  % \bibPrep is defined in header.tex
\bibliography{src/bibdb}

\newpage
\section{Grading Rubric (40 points)}

For programming problems, students must submit both the code and the output.

\subsection{Exercise 1 (11 points)}

\begin{enumerate}
\item For each label (Y or N), showing your work and using the correct formula gives 5 points
    (1 for each multiplicative term).
\item Showing how to use the scores (or probabilities) to obtain the correct prediction gives 1 point.
\end{enumerate}

\subsection{Exercise 2 (3 points)}

\begin{enumerate}
\item Read the dataset: 1 point
\item Compute the total delay time: 1 point
\item Assign Y or N label: 1 point
\end{enumerate}

\subsection{Exercise 3 (26 points)}

The official solution uses functions to avoid repetition of code.
(E.g., it has a function called `predict`, which takes
a flight's details (source, destination, airline) as input
and displays the posterior probabilities and the prediction.
This avoids repeating the code for each of the 4 inputs (a) to (d).)

The rubric below is written with the official solution in mind,
but students are free to follow a different outline.
Specifically, there is no penalty for code repetition;
points will simply be split across each repetition.

\begin{enumerate}
\item 11 points for code that computes appropriate aggregations:
    \begin{tightenum}
    \item 2 points to compute aggregations needed for $n$.
    \item 6 points to compute aggregations needed for $\widehat{n}$.
    \item 3 points to compute aggregations needed for $p$.
    \end{tightenum}
\item 9 points for code that uses aggregates to compute scores:
    \begin{tightenum}
    \item 2 points for the $P(C_k)$ term.
    \item 4 points for the correct numerator of the $P(x_i \mid C_k)$ term.
    \item 2 points for the correct denominator of $P(x_i \mid C_k)$ term.
    \item 1 point to combine these terms to get the scores for each label.
    \end{tightenum}
\item 2 points for code that uses the scores to compute posterior probabilities.
\item 4 points for outputs. For each of the 4 inputs (a) to (d):
    \begin{tightenum}
    \item 0.5 point for the correct posterior probabilities of Y and N.
    \item 0.5 point for the correct prediction (Delayed or Not delayed).
    \end{tightenum}
\end{enumerate}

\subsection{Extra credit}

\begin{enumerate}
\item 2 points for not hard-coding the number and names of input columns.
    This makes it easy to adapt the solution to different datasets.
\item For datasets with a very large number of rows, most operations take a very long time.
    Avoiding needless repetition can considerably speed up computation.
    2 points will be awarded to code that is close to optimal in terms of number of passes over data.
\end{enumerate}

\end{document}
